\section{Bounding the path-type sum}
\label{sec:path-sum}

We now bound the sum in \eqref{eq:moment-reduction}. The coincidence factor
at an interior $V$ layer can be large only when both neighboring layers
have the opposite type. Removing these isolated positions and bounding the
endpoint factors will leave a sum controlled by the number of runs of
$\mathcal D$ layers.

\begin{lemma}\label{lem:word-sum}
For all sufficiently large $n$, use the layer schedule
\eqref{eq:balanced-schedule} and the parameters $\alpha,\beta,\rho$ of
\eqref{eq:parameters}. For a word
$\tau\in\{\mathcal N,\mathcal D\}^n$, let $w(\tau)$ be its number of
switches, let $d_V(\tau),d_U(\tau)$ count its $\mathcal D$ positions in
the respective variable groups, and let $N_V(\tau)$ be its set of
$\mathcal N$ positions in group $V$. For $t\in N_V(\tau)$, let $h_t$ be
the number of internal vertices adjacent to layer $t$ whose two incident
layers both have type $\mathcal N$. There is an absolute constant $C$ such
that
\begin{equation}\label{eq:word-sum}
 \sum_{\tau\in\{\mathcal N,\mathcal D\}^n}
 n^{-w(\tau)}\alpha^{-d_V(\tau)}\beta^{d_U(\tau)}
 \prod_{t\in N_V(\tau)}\bigl(1+\alpha^{-1}n^{-h_t}\bigr)
 \le \exp(C\sqrt n).
\end{equation}
\end{lemma}

\begin{proof}
Denote the summand on the left of \eqref{eq:word-sum} by $W(\tau)$.
Call a position isolated if it is an interior $V$ layer with type
$\mathcal N$ and both its neighbors have type $\mathcal D$. Change every
isolated position to type $\mathcal D$, obtaining a word $F(\tau)$.
No two $V$ layers are adjacent, so none of their neighbors is changed.
It follows that $F(\tau)$ has no isolated position, and the values of
$h_t$ at its remaining $\mathcal N_V$ positions are unchanged.

Each changed position has $h_t=0$ before the change. Its two incident
switches disappear, its factor $1+\alpha^{-1}$ disappears, and a factor
$\alpha^{-1}$ is introduced. Switches removed at distinct positions are
distinct, even when those positions are separated by only one $U$ layer.
Thus, if $j$ positions were changed,
\begin{equation}\label{eq:isolated-flip}
 \frac{W(\tau)}{W(F(\tau))}
 =\left(\frac{1+\alpha}{n^2}\right)^j.
\end{equation}

The preimages can be counted exactly. For a word $y$ with no isolated
position, let $E(y)$ be the set of interior $V$ positions whose type and
both neighboring types are $\mathcal D$. Its preimages are obtained by
changing an arbitrary subset of $E(y)$ to $\mathcal N$. These positions
are nonadjacent, their neighbors remain unchanged, and the operation
creates precisely the isolated positions that are changed back by $F$.
Consequently
\[
 \sum_{\tau:F(\tau)=y}W(\tau)
 =W(y)\left(1+\frac{1+\alpha}{n^2}\right)^{|E(y)|}
 \le W(y)(1+2/n^2)^k,
\]
where $0<\alpha<1$ and $|E(y)|\le k$ were used.

In a word with no isolated position, every interior $\mathcal N_V$
position has $h_t\ge1$. An endpoint position can have $h_t=0$; there is
at most one endpoint $V$ layer, since the schedule starts with $V$ and
ends with $U$. Hence
\begin{equation}\label{eq:correction-removal}
 \sum_\tau W(\tau)
 \le (1+2/n^2)^k(1+\alpha^{-1})
       (1+\alpha^{-1}/n)^k Z,
 \qquad
 Z=\sum_\tau n^{-w(\tau)}
                 \alpha^{-d_V(\tau)}\beta^{d_U(\tau)}.
\end{equation}
Here we enlarged the final sum from words without isolated positions to
all words. By Lemma~\ref{lem:parameters}, $\alpha^{-1}\le\sqrt n$ and
$k\le n/2$, so the logarithm of the prefactor is at most
\[
 \frac{2k}{n^2}+\log(1+\sqrt n)+\frac{k}{\sqrt n}
 =O(\sqrt n).
\]
It remains to bound $Z$. A run of $\mathcal D$ positions is a maximal
nonempty interval of such positions. Each interior run has two
switches, contributing $n^{-2}$, whereas a run meeting exactly one
endpoint has one switch. We will show that the product of layer
weights on every run is at most $n$ times a rounding factor. Thus the
switch costs leave a factor $n^{-1}$ for every interior run.

The relation between $\alpha$ and $\beta$ reduces the layer weight
to a discrepancy between the two layer counts. For an interval $I$
of consecutive layers, define
\[
 d(I)=\#\{V\text{ layers in }I\}
       -\rho\,\#\{U\text{ layers in }I\}.
\]
We claim that
\begin{equation}\label{eq:interval-discrepancy}
 d(I)\le1+\rho
 \quad\text{for every interval }I,
 \qquad d([n])=0.
\end{equation}
Indeed, $r$ consecutive complete chunks of the schedule, starting after
chunk $j$, contain $r$ layers in $V$ and $r+E$ layers in $U$, where
\[
 E=\left\lfloor\frac{(j+r)s}{k}\right\rfloor
      -\left\lfloor\frac{js}{k}\right\rfloor,
 \qquad |E-rs/k|<1.
\]
Using $m=k+s$ and $\rho=k/m$, their discrepancy is
\[
 r-\rho(r+E)=\rho(rs/k-E)<\rho.
\]
An interval spanning several chunks consists of consecutive complete
chunks, possibly preceded by a proper terminal part of one chunk and
followed by a proper initial part of another. The terminal part contains
only $U$ layers and has nonpositive discrepancy; the initial part
contains at most one $V$ layer and has discrepancy at most $1$.
An interval contained in a single chunk also has discrepancy at most
$1$. This proves the first assertion of \eqref{eq:interval-discrepancy}.
The second is the identity $k-\rho m=0$.

By Lemma~\ref{lem:parameters}, the product of layer weights on a
$\mathcal D$ run $I$ is at most
\[
 \alpha^{-\#V(I)}\beta^{\#U(I)}
 \le \alpha^{-d(I)}\exp(C\#U(I)/b)
 \le n\exp(C\#U(I)/b).
\]
For the last inequality, the sign of $d(I)$ causes no difficulty:
\[
 \alpha^{-d(I)}\le\alpha^{-(1+\rho)}
 \le n^{(1+\rho)/2}\le n,
\]
since $0<\alpha<1$, $\alpha\ge n^{-1/2}$, and $\rho<1$.
The runs are disjoint, so all their rounding factors together contribute
at most $\exp(Cn/b)$.

Consider a mixed word, meaning one containing both types. Let $r$ be its
number of interior $\mathcal D$ runs and $e$ its number of
$\mathcal D$ runs meeting an endpoint. A mixed word has no run meeting
both endpoints, so $e\in\{0,1,2\}$ and
\[
 w(\tau)=2r+e.
\]
There are $r+e$ runs. Their weight, including the switch factors, is
therefore at most
\[
 n^{-(2r+e)}n^{r+e}\exp(Cn/b)
 =n^{-r}\exp(Cn/b).
\]
The initial type and the set of switch positions uniquely specify a
word. Thus at most $2\binom{n-1}{2r+e}$ words have given values of $r,e$.
Taking binomial coefficients outside their usual range to be zero,
the mixed-word contribution to $Z$ is at most
\[
 \exp(Cn/b)\sum_{e=0}^2\sum_{r\ge0}
             2\binom{n-1}{2r+e}n^{-r}
 \le 6n^2\exp(\sqrt n+Cn/b).
\]
To see the last inequality, for each $e\in\{0,1,2\}$ use
\[
 \sum_{r\ge0}\binom{n-1}{2r+e}n^{-r}
 \le \sum_{r\ge0}\frac{n^{r+e}}{(2r+e)!}
 \le n^2\sum_{r\ge0}\frac{(\sqrt n)^{2r}}{(2r)!}
 \le n^2\exp(\sqrt n).
\]
The constant word $\mathcal N^n$ has weight $1$. The other constant
word has weight
\[
 \alpha^{-k}\beta^m
 \le\alpha^{-k+\rho m}\exp(Cm/b)
 =\exp(Cm/b).
\]
It follows that
\[
 Z\le\exp(Cn/b)\bigl(2+6n^2\exp(\sqrt n)\bigr)
 =\exp(O(\sqrt n)),
\]
where $b=\Theta(n^{5/2})$. Combining this with
\eqref{eq:correction-removal} proves the lemma.
\end{proof}

\begin{proof}[Completion of the proof of Proposition~\ref{prop:imm}]
Lemma~\ref{lem:word-sum}, applied to \eqref{eq:moment-reduction}, gives
\[
 T_2\le D n^{2L}\mu^2(1+\beta)^m\exp(C\sqrt n)
      \le D n^{2L}\mu^2\exp(C'\sqrt n),
\]
because $\beta\le2/\sqrt n$ and $m\le n$.
By \eqref{eq:first-trace},
\[
 T\ge Dn^L\mu\exp(-C''n^{-1/2})>0.
\]
For the positive semidefinite matrix $H=F_f^*F_f$, the
Cauchy--Schwarz inequality applied to its nonzero eigenvalues yields
\[
 (\tr H)^2\le\rank(H)\,\tr(H^2).
\]
Since $T>0$ and $\rank(H)=\rank(F_f)=R(f)$, the two moment bounds imply
\[
 R(f)\ge\frac{T^2}{T_2}
       \ge D\exp(-C'''\sqrt n).
\]
This is the assertion of Proposition~\ref{prop:imm}.
\end{proof}
